\documentclass[11pt,professionalfont,aspectratio=43]{beamer}

% Assets
\usepackage[czech]{babel}				% Jazyk
\usepackage[a-2u]{pdfx}				    % Kopírování z pdfka
\usepackage{tikz}						% Schémata automatů
\usetikzlibrary{shadows, shadows.blur, shapes.geometric, arrows, positioning, overlay-beamer-styles, patterns, shadings, shapes.callouts}
% \usetikzlibrary{arrows}
\usepackage{pgfplots}
% \pgfplotsset{compat=1.15}
% \usepackage{mathrsfs}
\usepackage[export]{adjustbox}

\usepackage[T1]{fontenc}
\usefonttheme[onlymath]{serif}
% \usepackage{newtxmath}
\usepackage{pxfonts}
\usepackage{slantsc}
% \usepackage{varwidth}

\usepackage[linesnumbered,ruled,vlined]{algorithm2e}                % Pseudokód algoritmů
\usepackage[utf8]{inputenc}
% \usepackage{textcomp}
\usepackage{hyperref}
\usepackage{fancybox}

%%% FONTS
% \usefonttheme{serif}
% \usepackage{lmodern}
\usepackage{helvet}

\hypersetup{%
    pdfencoding=auto,
    pdfauthor={\insertauthor},
    pdftitle={\insertsubtitle}
}
\usepackage{float}
\usepackage{csquotes}					% české uvozovky
\usepackage{enumerate}					% enumerate environment
\usepackage{indentfirst}
\usepackage{tcolorbox}                  % Fancy boxíky
\usepackage{mathtools}
\usepackage{pifont}
\usepackage{cancel}
\usepackage{subcaption}
\usepackage{soul}
\usepackage{makecell}                   % Zalomení v buňce tabulky
\usepackage{xcolor}
\usepackage{graphicx}
\usepackage{tabularx}
\usepackage{amsmath}
\usepackage{amsthm}
\usepackage{colortbl}                   % Barvení buněk v tabulce
\usepackage{multicol}                   % Rozdělení obsahu do sloucpů

\usepackage{listings}                   % Úryvky z kódu
\definecolor{maroon}{rgb}{0.5,0,0}
\definecolor{forestgreen}{rgb}{0.13, 0.55, 0.13}

% \lstdefinelanguage{XML}
% {
%   basicstyle=\ttfamily,
%   morestring=[s]{"}{"},
%   morecomment=[s]{?}{?},
%   morecomment=[s]{!--}{--},
%   commentstyle=\color{forestgreen},
%   moredelim=[s][\color{black}]{>}{<},
%   moredelim=[s][\color{red}]{\ }{=},
%   stringstyle=\color{blue},
%   identifierstyle=\color{maroon}
% }
% Colors

\def\maincolR{0} % červená složka (0–255)
\def\maincolG{83} % zelená složka
\def\maincolB{120} % modrá složka
\definecolor{maincolor}{HTML}{8B00B5}

% \pgfmathtruncatemacro{\authorfooterbgR}{1*\maincolR}
% \pgfmathtruncatemacro{\authorfooterbgG}{1*\maincolG}
% \pgfmathtruncatemacro{\authorfooterbgB}{1*\maincolB}
% \definecolor{authorfooterbg}{RGB}{\authorfooterbgR,\authorfooterbgG,\authorfooterbgB}

% \pgfmathtruncatemacro{\titlefooterbgR}{0.722*\maincolR}   
% \pgfmathtruncatemacro{\titlefooterbgG}{0.974*\maincolG}   
% \pgfmathtruncatemacro{\titlefooterbgB}{0.739*\maincolB}
% \definecolor{titlefooterbg}{RGB}{\titlefooterbgR,\titlefooterbgG,\titlefooterbgB}

% Beamer theme
\colorlet{themecolor}{maincolor}
\colorlet{authorfooterbg}{maincolor}
\definecolor{titlefooterbg}{HTML}{9471B0}
\definecolor{titlefooterfg}{HTML}{000000}
\colorlet{titlecolorfg}{maincolor}
\definecolor{titleboxbg}{HTML}{E7E7E7}
\definecolor{datefooterbg}{HTML}{E3E3E3}
\colorlet{titleframeboxcolor}{maincolor}

% Background gradient
\definecolor{gradMiddle}{RGB}{212,248,255}
\definecolor{gradEnd}{RGB}{114,231,255}

% Block theme
\definecolor{blocktitlebg}{HTML}{F0DFA9}
\definecolor{blocktitlefg}{HTML}{957300}
\definecolor{blockbodybg}{HTML}{F5F0DF}

\definecolor{alertblockbodybg}{HTML}{FFE4E4}
\definecolor{alertblocktitlebg}{HTML}{ECA4A4}
\definecolor{alertblocktitlefg}{HTML}{B11313}

\definecolor{exampleblockbodybg}{HTML}{DAF0DA}
\definecolor{exampleblocktitlebg}{HTML}{BBE0BB}
\definecolor{exampleblocktitlefg}{HTML}{257A25}

\definecolor{mathfg}{HTML}{00185E}

\colorlet{bulletbg}{maincolor}
% \definecolor{bulletbg}{HTML}{}

% Beamer theme
\usetheme{Madrid}
\usecolortheme{seahorse}
% \useoutertheme[subsection=false]{miniframes}

\usecolortheme[named=themecolor]{structure}
\setbeamertemplate{frame numbering}[fraction]
\setbeamertemplate{navigation symbols}{}
\setbeamercovered{invisible}

% Header/footer
\setbeamercolor{author in head/foot}{bg=authorfooterbg,fg=white}
\setbeamercolor{title in head/foot}{bg=titlefooterbg,fg=titlefooterfg}
\setbeamercolor{date in head/foot}{bg=datefooterbg,fg=black}
% \setbeamercolor{title}{bg=datefooterbg}
% \setbeamercolor{title}{fg=titlecolorfg}
\setbeamerfont{title}{series=\bfseries}
\setbeamercolor{frametitle}{bg=datefooterbg}
% \setbeamercolor{section in head/foot}{bg=red,fg=cyan}
% \setbeamercolor{subsection in head/foot}{bg=blue,fg=yellow}

% Table of contents
\setbeamercolor{section in toc}{fg=black}
\setbeamercolor{subsection in toc}{fg=red}
\setbeamercolor{section number projected}{fg=black}

% Blocks
\setbeamertemplate{blocks}[rounded][shadow=true]

\setbeamercolor{block title}{bg=blocktitlebg, fg=blocktitlefg}
\setbeamercolor{block body}{parent=normal text,use=block title,bg=blockbodybg}

\setbeamercolor{block body alerted}{bg=alertblockbodybg}
\setbeamercolor{block title alerted}{bg=alertblocktitlebg,fg=alertblocktitlefg}

\setbeamercolor{block body example}{bg=exampleblockbodybg}
\setbeamercolor{block title example}{bg=exampleblocktitlebg,fg=exampleblocktitlefg}

% Math
\setbeamercolor{math text}{fg=mathfg}

% Itemize a enumerate
\setbeamertemplate{itemize items}[circle]
\setbeamertemplate{enumerate items}[circle]
\setbeamercolor{itemize item}{fg=bulletbg,bg=white}
\setbeamercolor{itemize subitem}{fg=bulletbg,bg=white}
\setbeamercolor{itemize subsubitem}{fg=bulletbg,bg=white}

\setbeamertemplate{enumerate items}[circle]
\setbeamercolor{item projected}{bg=bulletbg}

% Titlepage
\defbeamertemplate*{title page}{customized}[1][]
{
    \vbox{}%
    \vfill%
    \begin{centering}
        \begin{tikzpicture}
            \node[thick, rectangle, rounded corners=5mm, drop shadow={shadow xshift=0mm, shadow yshift=2mm, opacity=0.2, blur shadow, shadow blur steps=10}, inner sep=5mm, fill=titleboxbg, fill opacity=1] {%
                {\usebeamerfont{title}\textcolor{titlecolorfg}{\inserttitle}\par}%
            };
        \end{tikzpicture}
        \vskip1em\par%
        \ifx\insertsubtitle\@empty%
        \else%
            {\usebeamerfont{subtitle}\insertsubtitle\par}%
        \fi%
        \vskip1em\par%
        \ifx\insertauthor\@empty%
        \else%
            {\usebeamerfont{author}\insertauthor}\\%
            \texttt{\email}
        \fi
        \vskip0.3em\par%
        % \ifx\insertinstitute\@empty%
        % \else%
        %     {\usebeamerfont{institute}\insertinstitute\par}%
        % \fi%
        \vskip1em\par%
        \ifx\inserttitlegraphic\@empty%
        \else%
            \vskip1em\par%
            \inserttitlegraphic\par%
        \fi%
        \vskip1em\par%
        \ifx\insertdate\@empty%
        \else%
            {\usebeamerfont{date}\insertdate\par}%
        \fi%
    \end{centering}%
    \vfill%
}
% Enumerate
%\setlist[enumerate]{topsep=0pt,itemsep=-1ex,partopsep=1ex,parsep=1ex,label=(\arabic*)}

\MakeOuterQuote{"}

%%% Makra pro definice, věty, tvrzení, příklady, ... (vyžaduje baliček amsthm)

\theoremstyle{plain}
\newtheorem{thm}{Věta}[section]
\newtheorem{lem}[theorem]{Lemma}
\newtheorem{prop}[theorem]{Tvrzení}
\newtheorem{cor}[theorem]{Důsledek}
\newtheorem*{prop*}{Tvrzení}

\theoremstyle{definition}
\newtheorem{define}[theorem]{Definice}
% \newtheorem{example}[theorem]{Příklad}
% \newtheorem{remark}[theorem]{Poznámka}
% \newtheorem{convention}[theorem]{Úmluva}
% \newtheorem{denoting}[theorem]{Značení}

\newenvironment{defblock}[1][]{
  \begin{block}{\ifthenelse{\equal{#1}{}}
  {Definice}
  {Definice (#1)}}
}{\end{block}}
\newenvironment{thmblock}[1][]{
  \begin{block}{\ifthenelse{\equal{#1}{}}
  {Věta}
  {Věta (#1)}}
}{\end{block}}
\newenvironment{corblock}[1][]{
  \begin{block}{\ifthenelse{\equal{#1}{}}
  {Důsledek}
  {Důsledek (#1)}}
}{\end{block}}
\newenvironment{propblock}[1][]{
  \begin{block}{\ifthenelse{\equal{#1}{}}
  {Tvrzení}
  {Tvrzení (#1)}}
}{\end{block}}
\newenvironment{exblock}[1][]{
  \begin{exampleblock}{\ifthenelse{\equal{#1}{}}
  {Příklad}
  {Příklad (#1)}}
}{\end{exampleblock}}


% Colors
\definecolor{lightblue}{HTML}{009AD4}
\definecolor{lightgreen}{HTML}{68FF00}
\definecolor{darkgreen}{HTML}{00B600}
\definecolor{darkblue}{HTML}{0265b3}
\definecolor{darkred}{HTML}{DA0707}
\definecolor{lightred}{HTML}{FF5100}
\definecolor{orange}{HTML}{FFE000}
\definecolor{codeblue}{HTML}{007eed}
\definecolor{codegreen}{rgb}{0,0.6,0}
\definecolor{codegray}{rgb}{0.5,0.5,0.5}
\definecolor{codebeige}{HTML}{D4A000}
\definecolor{codepink}{HTML}{FF0055}
\definecolor{backcolour}{rgb}{0.95,0.95,0.92}

\definecolor{decproblemtitlefg}{HTML}{957300}
\definecolor{decproblembg}{HTML}{F5F0DF}

\definecolor{timportantboxcolor}{HTML}{A97100}
\definecolor{talertboxcolor}{HTML}{B70000}
\definecolor{tnoticeboxcolor}{HTML}{159C00}
\definecolor{tinfoboxcolor}{HTML}{004DFF}

\newcommand{\markred}[1]{\textcolor{lightred}{#1}}
\newcommand{\markdarkred}[1]{\textcolor{darkred}{#1}}
\newcommand{\markgreen}[1]{\textcolor{lightgreen}{#1}}
\newcommand{\markdarkgreen}[1]{\textcolor{darkgreen}{#1}}
\newcommand{\markorange}[1]{\textcolor{orange}{#1}}
\newcommand{\markblue}[1]{\textcolor{lightblue}{#1}}
\newcommand{\markdarkblue}[1]{\textcolor{darkblue}{#1}}

% Math mode settings
% \AtBeginDocument{
%     \setlength{\abovedisplayskip}{3pt}
%     \setlength{\belowdisplayskip}{-4pt}
% }

% Inline images
\newcommand{\inlineimgscale}{1.1}

% X and check mark
\newcommand{\cmark}{\markgreen{\ding{51}}}
\newcommand{\xmark}{\markred{\ding{55}}}

% Math
\newcommand{\R}{\mathbb{R}}
\newcommand{\C}{\mathbb{C}}
\newcommand{\N}{\mathbb{N}}
\newcommand{\Z}{\mathbb{Z}}
\newcommand{\Q}{\mathbb{Q}}

\newcommand{\T}[1]{#1^\top}                     % Tranpozice matice
\newcommand{\compl}[1]{#1^\complement}
\newcommand{\mapping}[3]{#1:#2 \rightarrow #3}
\newcommand{\mapto}[3]{#1:#2 \mapsto #3}
\newcommand{\set}[1]{\left\{#1\right\}}
\newcommand{\powset}{\mathcal{P}}
\newcommand{\code}[1]{\langle#1\rangle}
\DeclareMathOperator*{\argmin}{arg\,min}
\DeclareMathOperator*{\argmax}{arg\,max}

% Statistika a AI
\newcommand{\average}[1]{\overline{#1}}
\DeclareMathOperator{\median}{med}
\DeclareMathOperator{\Var}{var}
\DeclareMathOperator{\Cov}{cov}
\DeclareMathOperator{\MSE}{MSE}
\DeclareMathOperator{\SSE}{SSE}

\DeclareMathOperator{\TP}{TP}
\DeclareMathOperator{\TN}{TN}
\DeclareMathOperator{\FP}{FP}
\DeclareMathOperator{\FN}{FN}
\DeclareMathOperator{\Accuracy}{Acc}
\DeclareMathOperator{\Precision}{Prec}
\DeclareMathOperator{\Recall}{Rec}
\DeclareMathOperator{\Fscore}{F}
\DeclareMathOperator{\Gini}{Gini}

\DeclareMathOperator{\mspan}{span}
\DeclareMathOperator{\mrank}{rank}
\DeclareMathOperator{\tg}{tg}
\DeclareMathOperator{\id}{id}
\DeclareMathOperator{\dom}{dom}
\DeclareMathOperator{\rng}{rng}
\renewcommand{\emptyset}{\varnothing}
\newcommand{\binary}[1]{\ensuremath{\left\langle #1\right\rangle}_B}

% DFS in and out lists
\DeclareMathOperator{\inlist}{in}
\DeclareMathOperator{\outlist}{out}
% \DeclareMathOperator{\lowlist}{low}
\DeclareMathOperator{\key}{key}

% A* macro
\newcommand{\Astar}{$\textcolor{black}{\textup{A}^{*}}$}

% Complexity classes
\DeclareMathOperator{\classP}{P}
\DeclareMathOperator{\classNP}{NP}
\DeclareMathOperator{\classTime}{TIME}
\DeclareMathOperator{\classNTime}{NTIME}
\DeclareMathOperator{\classSpace}{SPACE}
\DeclareMathOperator{\classNSpace}{NSPACE}

\DeclareMathOperator{\SAT}{SAT}
\DeclareMathOperator{\THREESAT}{3-SAT}
\DeclareMathOperator{\UNSAT}{UNSAT}
\DeclareMathOperator{\NZMNA}{NzMna}
\DeclareMathOperator{\KLIKA}{Klika}
\DeclareMathOperator{\ROZVRH}{Rozvrh}
\DeclareMathOperator{\SUDOKU}{Sudoku}
\newcommand{\HALT}{\ensuremath{\textsc{Halt}}}
\newcommand{\FACTOR}{\ensuremath{\textsc{Factor}}}

\DeclareMathOperator{\regex}{RegE}

%%% Barevné boxíky (notice / info / important / alert)
%
% Šířka boxu se nezadává ručně -- box se přizpůsobí šířce svého obsahu.
% Obsah se vysází do prostředí varwidth, takže krátký text zůstane na
% jednom řádku; delší text se zalomí, nejvýše však na šířku \linewidth
% (tj. \textwidth mimo sloupce a minipage). Pokud je titulek širší než
% obsah, řídí šířku boxu titulek.

\newsavebox{\tboxcontentbox}
\newsavebox{\tboxtitlebox}
\newlength{\tboxwidth}
\newlength{\tboxpadding}

% Geometrie boxu (musí odpovídat stylu tautobox níže):
% 2 * (boxrule + boxsep + left)
\setlength{\tboxpadding}{\dimexpr 0.5mm+1mm+4mm\relax}
\setlength{\tboxpadding}{2\tboxpadding}

\tcbset{tautobox/.style={boxrule=0.5mm, boxsep=1mm, left=4mm, right=4mm,
                         fonttitle=\bfseries}}

% Začátek sběru obsahu do boxu
\newcommand{\tboxcollect}{%
  \begin{lrbox}{\tboxcontentbox}%
    \begin{varwidth}{\dimexpr\linewidth-\tboxpadding\relax}%
}

% Ukončení sběru a výpočet výsledné šířky
% #1 = titulek (pro měření; prázdný u variant bez titulku)
\newcommand{\tboxmeasure}[1]{%
    \end{varwidth}%
  \end{lrbox}%
  \sbox{\tboxtitlebox}{\bfseries #1}%
  \setlength{\tboxwidth}{\wd\tboxcontentbox}%
  \ifdim\wd\tboxtitlebox>\tboxwidth
    \setlength{\tboxwidth}{\wd\tboxtitlebox}%
  \fi
  \addtolength{\tboxwidth}{\tboxpadding}%
  \ifdim\tboxwidth>\linewidth
    \setlength{\tboxwidth}{\linewidth}%
  \fi
}

% Vysázení boxu s titulkem: #1 = barva, #2 = titulek
\newcommand{\tboxrenderbox}[2]{%
  \tboxmeasure{#2}%
  \begin{center}%
    \begin{tcolorbox}[tautobox, colback=#1!15, colframe=#1,
                      title={#2}, width=\tboxwidth]
      \usebox{\tboxcontentbox}%
    \end{tcolorbox}%
  \end{center}%
}

% Vysázení boxu bez titulku: #1 = barva
\newcommand{\tboxrenderframe}[1]{%
  \tboxmeasure{}%
  \begin{center}%
    \begin{tcolorbox}[tautobox, colback=#1!15, colframe=#1,
                      width=\tboxwidth]
      \usebox{\tboxcontentbox}%
    \end{tcolorbox}%
  \end{center}%
}

% Poznámky (notice)
\newenvironment{tnoticebox}[1]{%
  \def\tboxtitletext{#1}\tboxcollect
}{%
  \tboxrenderbox{tnoticeboxcolor}{\tboxtitletext}%
}

\newenvironment{tnoticeframe}{%
  \tboxcollect
}{%
  \tboxrenderframe{tnoticeboxcolor}%
}

% Informační bloky (info)
\newenvironment{tinfobox}[1]{%
  \def\tboxtitletext{#1}\tboxcollect
}{%
  \tboxrenderbox{tinfoboxcolor}{\tboxtitletext}%
}

\newenvironment{tinfoframe}{%
  \tboxcollect
}{%
  \tboxrenderframe{tinfoboxcolor}%
}

% Důležité bloky (important)
\newenvironment{timportantbox}[1]{%
  \def\tboxtitletext{#1}\tboxcollect
}{%
  \tboxrenderbox{timportantboxcolor}{\tboxtitletext}%
}

\newenvironment{timportantframe}{%
  \tboxcollect
}{%
  \tboxrenderframe{timportantboxcolor}%
}

% Varovné bloky (alert)
\newenvironment{talertbox}[1]{%
  \def\tboxtitletext{#1}\tboxcollect
}{%
  \tboxrenderbox{talertboxcolor}{\tboxtitletext}%
}

\newenvironment{talertframe}{%
  \tboxcollect
}{%
  \tboxrenderframe{talertboxcolor}%
}

% Decision problem
\newcommand{\decproblem}[3]{
    \begin{center}
        \begin{minipage}{.9\textwidth}
            \begin{table}[!htbp]
                \centering
                \bgroup
                \def\arraystretch{1.8}
                \begin{tabularx}{\textwidth}{|rX|}
                \hline
                \rowcolor{decproblembg} \multicolumn{2}{|c|}{\textcolor{decproblemtitlefg}{\textsc{#1}}} \\ \hline
                \markdarkred{\textbf{Vstup:}} & #2                \\ \hline
                \rowcolor{decproblembg} \markdarkred{\textbf{Otázka:}} & #3               \\ \hline
                \end{tabularx}
                \egroup
            \end{table}
        \end{minipage}
    \end{center}
}
\newcommand{\bigO}{\mathcal{O}}

% TODO
\newcommand{\todo}[1]{\textcolor{red}{(\noindent TODO: #1.)}}


\newcommand{\hlcode}[1]{\colorbox{red}{#1}}

% Algorithms
\numberwithin{algocf}{section}
\renewcommand{\algorithmcfname}{Algoritmus}
\SetKwInput{KwIn}{Vstup}
\SetKwInput{KwOut}{Výstup}
\SetKwProg{Function}{Funkce}{}{end}
\setlength{\algomargin}{25pt}
\renewcommand{\thealgocf}{}     % Vypnutí číslování algoritmů

% Definice stylů pro uzly
\tikzstyle{startstop} = [rectangle, rounded corners, minimum width=1.5cm, minimum height=0.5cm,text centered, draw=black, fill=red!30]
\tikzstyle{process}   = [rectangle, minimum width=1.5cm, minimum height=0.5cm, text centered, draw=black, fill=orange!30]
\tikzstyle{decision}  = [diamond, aspect=2, text centered, draw=black, fill=green!30]
\tikzstyle{arrow}     = [thick,->,>=stealth]

\tikzstyle{vertex}    = [thick, draw, circle, minimum size=4mm, inner sep=1pt, outer sep=1pt, align=center]
\tikzstyle{vertexplain}    = [minimum size=4mm, inner sep=1pt, outer sep=1pt, align=center]
\tikzstyle{verteximg}    = [draw, circle, minimum size=4mm, inner sep=0pt, outer sep=1pt, align=center, clip]
\tikzstyle{verteximgplain}    = [minimum size=4mm, inner sep=0pt, outer sep=1pt, align=center]
\tikzstyle{task}    = [draw, rounded corners=2pt, minimum width=14mm, minimum height=8mm, inner sep=2pt, align=center]

\tikzstyle{edge}    = [thick,>=stealth]
\tikzstyle{diredge}    = [thick, ->, >=stealth]

%% Automaty
\tikzstyle{state}=[%
    draw,
    circle,
    thick,
    minimum size=10mm,
    inner sep=0pt,
    outer sep=0pt
]

\tikzstyle{acceptstate}=[%
    state,
    double,
    double distance=1pt
]

\tikzstyle{transition}=[%
    ->,
    >=stealth,
    thick
]

\newcommand<>{\state}[4][]{%
    \node#5[state,#1] (#2) at #3 {#4};
}

\newcommand<>{\acceptingstate}[4][]{%
    \node#5[acceptstate,#1] (#2) at #3 {#4};
}

\newcommand<>{\transitionarrow}[5][]{%
    \path#6 (#2) edge[transition,#1] node[#3] {#4} (#5);
}

\newcommand<>{\initialstate}[6][]{%
    \node#7[state,#1] (#2) at #3 {#4};
    \draw#7[transition] #5 -- (#2.#6);
}

%%%%%

% Frames
\newcommand{\tableofcontentsframe}{
    \begin{frame}{Obsah}
        \tableofcontents
    \end{frame}
}

% \definecolor{boxcolor}{HTML}{00820a}
\newcommand{\titleframe}[1]{%
  \begin{frame}[plain]
    \begin{tikzpicture}[remember picture, overlay]
      \node at (current page.center)
        [draw=titleframeboxcolor, line width=1pt,
         inner xsep=10mm, inner ysep=4mm, rounded corners=2mm,
         text width=0.8\textwidth, align=center,
         execute at begin node=\setlength{\baselineskip}{2.2em}]
        {\Huge\markred{#1}\normalfont};
    \end{tikzpicture}
  \end{frame}
  \addtocounter{framenumber}{-1}%
}

\newenvironment{titleframeenv}[1]{
    \begin{frame}[plain]
        \begin{tcolorbox}[colback=white,boxrule=1pt]
            \begin{center}
                \Huge\markred{#1}\normalfont
            \end{center}
        \end{tcolorbox}
}{\addtocounter{framenumber}{-1}\end{frame}}

\newcommand{\icon}[1]{
  \includegraphics[height=\fontcharht\font`\B]{#1}
}

% Obrázek, který se vždy vejde na snímek. Adjustbox jej v případě potřeby
% zmenší tak, aby nepřesáhl šířku textu ani zadanou část výšky snímku; díky
% tomu sedí sazba ve všech používaných poměrech stran (4:3, 16:9 i 16:10).
% Volitelným argumentem je maximální výška (výchozí 0.68\textheight), povinným
% cesta k souboru s obrázkem (TikZ zdroj vkládaný přes \input).
\newcommand{\obrazek}[2][0.68\textheight]{%
    \begin{figure}
        \centering
        \begin{adjustbox}{max width=\linewidth, max totalheight=#1, center}
            \input{#2}
        \end{adjustbox}
    \end{figure}%
}

% Totéž pro rastrové či vektorové obrázky vkládané přes \includegraphics.
\newcommand{\obrazekgrf}[2][0.68\textheight]{%
    \begin{figure}
        \centering
        \includegraphics[width=\linewidth, max totalheight=#1, center]{#2}
    \end{figure}%
}

% Závěrečný snímek se seznamem zdrojů. Prvním (volitelným) argumentem je název
% sekce i snímku, druhým čárkou oddělený seznam klíčů z ../assets/literatura.bib.
% Citace se sázejí podle ČSN ISO 690, v pořadí, v jakém jsou klíče uvedeny.
\newcommand{\zdrojeframe}[2][Zdroje]{%
    \section{#1}
    \nocite{#2}%
    \begin{frame}[allowframebreaks]{#1}
        \printbibliography[heading=none]
    \end{frame}
}
\input{../assets/listing.tex}

\def\presentationtitle{Dijsktrův algoritmus}

\definecolor{dijunknown}{RGB}{190,190,190}
\definecolor{dijopen}{RGB}{255,230,120}
\definecolor{dijclosed}{RGB}{170,220,170}
\definecolor{dijtree}{RGB}{255,0,0}

\tikzset{
    dijvertex/.style={
        circle,
        draw=black,
        thick,
        minimum size=10mm,
        inner sep=0pt,
        font=\scriptsize,
        align=center
    },
    dijedge/.style={
        thick
    },
    dijtreeedge/.style={
        very thick,
        draw=dijtree
    },
    dijbox/.style={
        draw=black,
        rounded corners=2mm,
        fill=white,
        inner sep=2mm,
        font=\small
    },
    dijweight/.style={
        fill=white,
        inner sep=1pt,
        font=\scriptsize
    }
}

% vrchol s popiskem a aktuální vzdáleností
% #1 = jméno souřadnice (v1,...,v10)
% #2 = barva výplně
% #3 = popisek vrcholu
% #4 = aktuální odhad vzdálenosti
\newcommand{\dijnode}[4]{%
    \node[dijvertex, fill=#2] at (#1) {#3\\[1.2mm]#4};
}

% #1 = titulek kroku
% #2 = obsah priority queue
% #3 = zvýrazněné hrany stromu nejkratších cest
% #4 = vrcholy
\newcommand{\DijkstraState}[4]{%
\begin{tikzpicture}[scale=0.8,x=1.55cm,y=1.2cm]

    % --- souřadnice vrcholů ---
    \coordinate (v1)  at (0,0);
    \coordinate (v2)  at (1.9,1.6);
    \coordinate (v3)  at (1.9,-1.8);
    \coordinate (v4)  at (4.0,2.4);
    \coordinate (v5)  at (4.0,0.8);
    \coordinate (v6)  at (4.0,-0.9);
    \coordinate (v7)  at (4.0,-2.6);
    \coordinate (v8)  at (6.3,1.6);
    \coordinate (v9)  at (6.3,-0.2);
    \coordinate (v10) at (8.5,0.8);

    % --- všechny hrany grafu ---
    \draw[dijedge] (v1)-- node[dijweight] {$4$} (v2);
    \draw[dijedge] (v1)-- node[dijweight] {$2$} (v3);
    \draw[dijedge] (v1)-- node[dijweight] {$7$} (v5);

    \draw[dijedge] (v2)-- node[dijweight] {$1$} (v3);
    \draw[dijedge] (v2)-- node[dijweight] {$5$} (v4);
    \draw[dijedge] (v2)-- node[dijweight] {$2$} (v5);
    \draw[dijedge] (v2)-- node[dijweight] {$6$} (v6);

    \draw[dijedge] (v3)-- node[dijweight] {$3$} (v5);
    \draw[dijedge] (v3)-- node[dijweight] {$8$} (v6);
    \draw[dijedge] (v3)-- node[dijweight] {$10$} (v7);

    \draw[dijedge] (v4)-- node[dijweight] {$2$} (v5);
    \draw[dijedge] (v4)-- node[dijweight] {$4$} (v8);

    \draw[dijedge] (v5)-- node[dijweight] {$2$} (v6);
    \draw[dijedge] (v5)-- node[dijweight] {$5$} (v8);
    \draw[dijedge] (v5)-- node[dijweight] {$6$} (v9);

    \draw[dijedge] (v6)-- node[dijweight] {$3$} (v7);
    \draw[dijedge] (v6)-- node[dijweight] {$2$} (v9);

    \draw[dijedge] (v7)-- node[dijweight] {$4$} (v9);

    \draw[dijedge] (v8)-- node[dijweight] {$1$} (v9);
    \draw[dijedge] (v8)-- node[dijweight] {$3$} (v10);

    \draw[dijedge] (v9)-- node[dijweight] {$2$} (v10);

    % --- zvýrazněné hrany stromu nejkratších cest ---
    #3

    % --- vrcholy ---
    #4

    % --- titulek kroku ---
    \node[anchor=west, font=\bfseries] at (-0.2,3.3) {#1};

    % --- legenda ---
    % \node[
    %     dijbox,
    %     anchor=north west,
    %     minimum width=4.8cm,
    %     minimum height=3.7cm
    % ] (legend) at (9.4,3.0) {};

    % \node[dijvertex, fill=dijunknown, minimum size=4.5mm]
    %     (leg1) at ([xshift=7mm,yshift=-7mm]legend.north west) {};
    % \node[anchor=west] at ([xshift=8mm]leg1.east) {nezpracovaný};

    % \node[dijvertex, fill=dijopen, minimum size=4.5mm]
    %     (leg2) at ([yshift=-4mm]leg1.south) {};
    % \node[anchor=west] at ([xshift=8mm]leg2.east) {odhadnutý};

    % \node[dijvertex, fill=dijclosed, minimum size=4.5mm]
    %     (leg3) at ([yshift=-4mm]leg2.south) {};
    % \node[anchor=west] at ([xshift=8mm]leg3.east) {uzavřený};

    % \draw[dijtreeedge] ([xshift=-2mm,yshift=-8mm]leg3.south west) -- ++(8mm,0);
    % \node[anchor=west] at ([xshift=10mm,yshift=-8mm]leg3.south west)
    %     {strom nejkratších cest};

    % % --- priority queue ---
    % \node[dijbox, anchor=north west, text width=4.8cm, align=left] at (9.4,-0.5) {%
    %     \textbf{Priority queue:}\\[1mm]
    %     #2
    % };

\end{tikzpicture}%
}

\begin{document}

    \title[\presentationtitle]{\textsc{Teoretická informatika}}
\subtitle{\presentationtitle}
\author{David Weber}
\def\email{weber3@spsejecna.cz}
\institute{SPŠE Ječná}
\date{%
    Rok 
    \pgfmathsetmacro{\firstyear}{ifthenelse(\month<9, \year-1, \year)}%
    \pgfmathsetmacro{\secondyear}{ifthenelse(\month<9, \year, \year+1)}%
    \pgfmathtruncatemacro{\secondyearlasttwo}{mod(\secondyear,100)}%
    \pgfmathprintnumber[fixed,precision=0,zerofill,1000 sep={}]{\firstyear}/%
    \pgfmathprintnumber[fixed,precision=0,zerofill,1000 sep={}]{\secondyearlasttwo}%
}
\titlegraphic{\includegraphics[width=0.25\textwidth]{../images/SPSE-Jecna_Logo.pdf}}

\begin{frame}[plain]
    \titlepage
\end{frame}

    \section{Opakování}

    \begin{frame}{Co již známe\dots}
        \begin{itemize}
            \item Již známe algoritmy BFS a DFS.
            \item Algoritmus BFS nalezne vždy nejkratší cestu do všech \textbf{dosažitelných vrcholů} v \textbf{neohodnoceném grafu}.
            \item Časová i prostorová složitost činí $\bigO(n+m)$.
        \end{itemize}
        \begin{tnoticeframe}
            Co když budeme hledat cestu v \textbf{ohodnoceném grafu}?
        \end{tnoticeframe}
    \end{frame}

    \section{Ohodnocené grafy}

    \titleframe{Ohodnocené grafy}

    \subsection{Definice ohodnoceného grafu}

    \begin{frame}{Ohodnocený graf}
        \begin{defblock}[Ohodnocený graf]
            \markdarkred{Ohodnocený grafem} rozumíme uspořádanou trojici $(V,E,w)$, přičemž $G=(V,E)$ je libovolný graf a $\mapping{w}{E}{\R}$ je zobrazení přiřazující každé hraně $e\in E$ její \markdarkred{váhu} $w(e)$.
        \end{defblock}
        \begin{tinfobox}[\textwidth]{Poznámka}
            Někdy je výhodnější se odkázat na samotné koncepvé vrcholy hrany $u,v$. Abychom si usnadnili zápis, budeme psát $w(u,v)$ místo $w(\set{u,v})$ v případě \textbf{neorientovaného grafu}, resp. místo $w((u,v))$ v případě \textbf{orientovaného grafu}.
        \end{tinfobox}
    \end{frame}

    \begin{frame}{Ohodnocený graf}{Příklad}
        \begin{figure}
            \centering
            \begin{tikzpicture}[scale=1.1]
    % graph
    \node[vertex] (a) at (0,0) {$v_1$};
    \node[vertex] (b) at (2,1.5) {$v_2$};
    \node[vertex] (c) at (4,0) {$v_3$};
    \node[vertex] (d) at (1.2,-2) {$v_4$};
    \node[vertex] (e) at (3.6,-2) {$v_5$};

    \draw[edge] (a) -- node[above left] {$2$} (b);
    \draw[edge] (b) -- node[above right] {$4$} (c);
    \draw[edge] (a) -- node[left] {$5$} (d);
    \draw[edge] (b) -- node[left] {$1$} (d);
    \draw[edge] (b) -- node[right] {$3$} (e);
    \draw[edge] (c) -- node[right] {$6$} (e);
    \draw[edge] (d) -- node[below] {$2$} (e);

    % mapping description
    \node[anchor=west, align=left] at (5.3,0.0) {%
        $\begin{aligned}
        w(\set{v_1,v_2}) &= w(v_1,v_2) = 2,\\
        w(\set{v_2,v_3}) &= w(v_2,v_3) = 4,\\
        w(\set{v_1,v_4}) &= w(v_1,v_4) = 5,\\
        w(\set{v_2,v_4}) &= w(v_2,v_4) = 1,\\
        w(\set{v_2,v_5}) &= w(v_2,v_5) = 3,\\
        w(\set{v_3,v_5}) &= w(v_3,v_5) = 6,\\
        w(\set{v_4,v_5}) &= w(v_4,v_5) = 2
        \end{aligned}$
    };
\end{tikzpicture}
        \end{figure}
    \end{frame}

    \subsection{Vzdálenost vrcholů v ohodnoceném grafu}

    \begin{frame}{Vzdálenost vrcholů v ohodnoceném grafu}
        \begin{defblock}
            Nechť $G=(V,E,w)$ je ohodnocený graf, vrcholy $u,v\in V$. Pak vzdáleností vrcholů $u,v$ rozumíme hodnotu
            \[d(u,v)=\min\set{\sum_{e\in P,e\in E}w(e):\text{$P$ je cesta z $u$ do $v$}}.\]
            Pokud $v$ není z $u$ dosažitelný, pak definujeme $d(u,v)=\infty$.
        \end{defblock}
        \begin{figure}
            \centering
            \begin{tikzpicture}[
    scale=1,
    every node/.style={font=\small},
    vertex/.style={draw, circle, thick, minimum size=7mm, inner sep=0pt},
    edge/.style={thick},
    selected/.style={draw=red!70!black, very thick},
    weight/.style={fill=white, inner sep=1.5pt}
]

% vrcholy
\node[vertex] (A) at (0,1.5) {$A$};
\node[vertex] (B) at (2,2.5) {$B$};
\node[vertex] (C) at (2,0.5) {$C$};
\node[vertex] (D) at (4,2.5) {$D$};
\node[vertex] (E) at (4,0.5) {$E$};
\node[vertex] (F) at (6,2.5) {$F$};
\node[vertex] (G) at (6,0.5) {$G$};
\node[vertex] (H) at (8,1.5) {$H$};

% běžné hrany
\draw[edge] (A) -- node[weight] {$3$} (C);
\draw[edge] (B) -- node[weight] {$2$} (C);
\draw[edge] (C) -- node[weight] {$5$} (E);
\draw[edge] (C) -- node[weight] {$6$} (D);
\draw[edge] (D) -- node[weight] {$3$} (E);
\draw[edge] (E) -- node[weight] {$4$} (G);
\draw[edge] (E) -- node[weight] {$5$} (F);
\draw[edge] (F) -- node[weight] {$4$} (G);
\draw[edge] (G) -- node[weight] {$5$} (H);

% vyznačená nejkratší cesta A-B-D-F-H
\draw[selected] (A) -- node[weight] {$2$} (B);
\draw[selected] (B) -- node[weight] {$3$} (D);
\draw[selected] (D) -- node[weight] {$2$} (F);
\draw[selected] (F) -- node[weight] {$3$} (H);

% popisek vzdálenosti
\node[align=center] at (4,-1.2) {\large $\displaystyle d(A,H)=2+3+2+3=10$};

\end{tikzpicture}
        \end{figure}
    \end{frame}

    \subsection{Hledání cest v ohodnocených grafech}

    \begin{frame}{Hledání cest v ohodnocených grafech}
        \visible<2->{\begin{figure}
            \centering
            \input{images/ohodnoceny-graf-nejkratsi-cesta.tex}
        \end{figure}}
        \visible<3->{\begin{tinfoframe}[0.6\textwidth]
            \begin{center}
                Jaká je nejkratší cesta z $u_1$ do $u_4$?
            \end{center}
        \end{tinfoframe}}
        \visible<4->{\begin{itemize}
            \item BFS by nalezl cestu $(u_1,u_3,u_4)$.
            \item Nicméně nejkratší cesta je $(u_1,u_2,u_3,u_4)$, \textbf{přestože má více hran!}
        \end{itemize}}
        \visible<5->{\begin{figure}
            \centering
            \begin{tikzpicture}[scale=1]
    
    \node[vertex] (U1) at (0.0,0.0) {$u_1$};
    \node[vertex] (U2) at (1.5,1.5) {$u_2$};
    \node[vertex] (U3) at (3.0,0.0) {$u_3$};
    \node[vertex] (U4) at (4.5,0.0) {$u_4$};

    \draw[edge,color=red] (U1) -- node[above,midway] {$1$} (U2);
    \draw[edge,color=red] (U2) -- node[above,midway] {$1$} (U3);
    \draw[edge,color=red] (U3) -- node[above,midway] {$4$} (U4);
    \draw[edge] (U1) -- node[above,midway] {$3$} (U3);

\end{tikzpicture}
        \end{figure}}
    \end{frame}

    \section{Dijkstrův algoritmus}

    \titleframe{Dijkstrův algoritmus}

    \subsection{Myšlenka a postup Dijkstrova algoritmu}

    \begin{frame}{Dijkstrův algoritmus}
        \begin{itemize}
            \visible<2->{\item Egster Dijkstra (1958)}
            \visible<3->{\item Na vstupu máme \textbf{nezáporně} ohodnocený graf $G=(V,E,w)$ a počáteční vrchol $v_0\in V$.}
            \visible<4->{\begin{tinfoframe}[\textwidth]
                \begin{center}
                    Zatím není zcela jasné, proč chceme, aby váhy hran byly nezáporné, ale později si to zdůvodníme.
                \end{center}
            \end{tinfoframe}}
            \visible<5->{\item Stejně jako v případě BFS i zde bychom rádi \textbf{při uzavírání vrcholu} chtěli znát nejkratší cestu do něj.}
            \visible<6->{\item Nicméně již nemusí platit, že \textbf{první nalezená cesta} od libovolného vrcholu $v$ je nutně nejkratší.}
            \visible<7->{\begin{timportantframe}[0.7\textwidth]
                \begin{center}
                    Cestu do $v$ je třeba postupně \textbf{zlepšovat}!
                \end{center}
            \end{timportantframe}}
            \visible<8->{\item Stejně jako v případě BFS si budeme uchovávat \markdarkred{stavy} a \markdarkred{vzdáleností} vrcholů.}
        \end{itemize}
    \end{frame}

    \begin{frame}{Postup Dijkstrova algoritmu}
        \begin{enumerate}
            \visible<2->{\item Inicializujeme seznamy stavů a vzdáleností vrcholů (vrcholy budou z počátku nenalezené a jejich vzdálenosti $\infty$).}
            \visible<3->{\item Označíme počáteční vrchol $v_0$ jako \textbf{otevřený} a nastavíme $D(v_0)=0$.}
            \visible<4->{\item Vybereme otevřený vrchol (pokud nějaký existuje) $v$ s nejmenší hodnotou $D(v)$.}
            \visible<5->{\item Vezmeme všechny sousedy $s$ vrcholu $v$ a pokud z některého vede kratší cesta do $v$, \textbf{aktualizujeme $D(v)$} a $s$ označíme jako \textbf{otevřený}.}
            \visible<6->{\item \textbf{Uzavřeme} vrchol $v$, vracíme ke 3. kroku a postup opakujeme.}
        \end{enumerate}
    \end{frame}

    \begin{frame}{Pseudokód Dijkstrova algoritmu}
        \begin{center}
            \begin{minipage}{0.9\textwidth}
                \begin{algorithm}[H]
                    \caption{\textsc{Dijkstra}}
                    \KwIn{Nezáporně ohodnocený graf $G=(V,E,w)$ a~počáteční vrchol $v_0\in V$}
                    \ForEach{\textup{vrchol $v\in V$}}{
                        $stav(v)\gets\textit{nenalezený},D(v)\gets\infty$\\
                    }
                    $stav(v_0)\gets\textit{otevřený},D(v_0)\gets 0$\\
                    \While{\textup{existují otevřené vrcholy}}{
                        % \tcp{Kandidát na nejkraší cestu do sousedních vrcholů}
                        $v\gets$ otevřený vrchol s~nejmenším $D$\\
                        
                        \ForEach{\textup{soused $s\in V$ vrcholu $v$}}{
                            % \tcp{Našli jsme lepší cestu}
                            \If{$D(v)+w(v,s)<D(s)$}{
                                $D(s)\gets D(v)+w(v,s)$\\
                                $stav(s)\gets\textit{otevřený}$
                            }
                        }
                        $stav(v)\gets\textit{uzavřený}$
                    }
                    \KwOut{Seznam vzdáleností $D$}
                \end{algorithm}
            \end{minipage}
        \end{center}
    \end{frame}

    \subsection{Dijkstrův algoritmus - příklad}

    \begin{frame}{Dijkstrův algoritmus}{Příklad}
    \centering
    \DijkstraState
    {Krok 0: inicializace}
    {$v_1:0$}
    {}
    {
        \dijnode{v1}{dijopen}{$v_1$}{$0$}
        \dijnode{v2}{dijunknown}{$v_2$}{$\infty$}
        \dijnode{v3}{dijunknown}{$v_3$}{$\infty$}
        \dijnode{v4}{dijunknown}{$v_4$}{$\infty$}
        \dijnode{v5}{dijunknown}{$v_5$}{$\infty$}
        \dijnode{v6}{dijunknown}{$v_6$}{$\infty$}
        \dijnode{v7}{dijunknown}{$v_7$}{$\infty$}
        \dijnode{v8}{dijunknown}{$v_8$}{$\infty$}
        \dijnode{v9}{dijunknown}{$v_9$}{$\infty$}
        \dijnode{v10}{dijunknown}{$v_{10}$}{$\infty$}
    }
    \end{frame}

    \begin{frame}{Dijkstrův algoritmus}{Příklad}
    \centering
    \DijkstraState
    {Krok 1: uzavíráme $v_1$}
    {$v_3:2,\ v_2:4,\ v_5:7$}
    {}
    {
        \dijnode{v1}{dijclosed}{$v_1$}{$0$}
        \dijnode{v2}{dijopen}{$v_2$}{$4$}
        \dijnode{v3}{dijopen}{$v_3$}{$2$}
        \dijnode{v4}{dijunknown}{$v_4$}{$\infty$}
        \dijnode{v5}{dijopen}{$v_5$}{$7$}
        \dijnode{v6}{dijunknown}{$v_6$}{$\infty$}
        \dijnode{v7}{dijunknown}{$v_7$}{$\infty$}
        \dijnode{v8}{dijunknown}{$v_8$}{$\infty$}
        \dijnode{v9}{dijunknown}{$v_9$}{$\infty$}
        \dijnode{v10}{dijunknown}{$v_{10}$}{$\infty$}
    }
    \end{frame}

    \begin{frame}{Dijkstrův algoritmus}{Příklad}
    \centering
    \DijkstraState
    {Krok 2: uzavíráme $v_3$}
    {$v_2:3,\ v_5:5,\ v_6:10,\ v_7:12$}
    {
        \draw[dijtreeedge] (v1)--(v3);
    }
    {
        \dijnode{v1}{dijclosed}{$v_1$}{$0$}
        \dijnode{v2}{dijopen}{$v_2$}{$3$}
        \dijnode{v3}{dijclosed}{$v_3$}{$2$}
        \dijnode{v4}{dijunknown}{$v_4$}{$\infty$}
        \dijnode{v5}{dijopen}{$v_5$}{$5$}
        \dijnode{v6}{dijopen}{$v_6$}{$10$}
        \dijnode{v7}{dijopen}{$v_7$}{$12$}
        \dijnode{v8}{dijunknown}{$v_8$}{$\infty$}
        \dijnode{v9}{dijunknown}{$v_9$}{$\infty$}
        \dijnode{v10}{dijunknown}{$v_{10}$}{$\infty$}
    }
    \end{frame}

    \begin{frame}{Dijkstrův algoritmus}{Příklad}
    \centering
    \DijkstraState
    {Krok 3: uzavíráme $v_2$}
    {$v_5:5,\ v_4:8,\ v_6:9,\ v_7:12$}
    {
        \draw[dijtreeedge] (v1)--(v3);
        \draw[dijtreeedge] (v3)--(v2);
    }
    {
        \dijnode{v1}{dijclosed}{$v_1$}{$0$}
        \dijnode{v2}{dijclosed}{$v_2$}{$3$}
        \dijnode{v3}{dijclosed}{$v_3$}{$2$}
        \dijnode{v4}{dijopen}{$v_4$}{$8$}
        \dijnode{v5}{dijopen}{$v_5$}{$5$}
        \dijnode{v6}{dijopen}{$v_6$}{$9$}
        \dijnode{v7}{dijopen}{$v_7$}{$12$}
        \dijnode{v8}{dijunknown}{$v_8$}{$\infty$}
        \dijnode{v9}{dijunknown}{$v_9$}{$\infty$}
        \dijnode{v10}{dijunknown}{$v_{10}$}{$\infty$}
    }
    \end{frame}

    \begin{frame}{Dijkstrův algoritmus}{Příklad}
    \centering
    \DijkstraState
    {Krok 4: uzavíráme $v_5$}
    {$v_4:7,\ v_6:7,\ v_8:10,\ v_9:11,\ v_7:12,\ v_{10}:17$}
    {
        \draw[dijtreeedge] (v1)--(v3);
        \draw[dijtreeedge] (v3)--(v2);
        \draw[dijtreeedge] (v3)--(v5);
    }
    {
        \dijnode{v1}{dijclosed}{$v_1$}{$0$}
        \dijnode{v2}{dijclosed}{$v_2$}{$3$}
        \dijnode{v3}{dijclosed}{$v_3$}{$2$}
        \dijnode{v4}{dijopen}{$v_4$}{$7$}
        \dijnode{v5}{dijclosed}{$v_5$}{$5$}
        \dijnode{v6}{dijopen}{$v_6$}{$7$}
        \dijnode{v7}{dijopen}{$v_7$}{$12$}
        \dijnode{v8}{dijopen}{$v_8$}{$10$}
        \dijnode{v9}{dijopen}{$v_9$}{$11$}
        \dijnode{v10}{dijopen}{$v_{10}$}{$17$}
    }
    \end{frame}

    \begin{frame}{Dijkstrův algoritmus}{Příklad}
    \centering
    \DijkstraState
    {Krok 5: uzavíráme $v_4$}
    {$v_6:7,\ v_8:10,\ v_9:11,\ v_7:12,\ v_{10}:17$}
    {
        \draw[dijtreeedge] (v1)--(v3);
        \draw[dijtreeedge] (v3)--(v2);
        \draw[dijtreeedge] (v3)--(v5);
        \draw[dijtreeedge] (v5)--(v4);
    }
    {
        \dijnode{v1}{dijclosed}{$v_1$}{$0$}
        \dijnode{v2}{dijclosed}{$v_2$}{$3$}
        \dijnode{v3}{dijclosed}{$v_3$}{$2$}
        \dijnode{v4}{dijclosed}{$v_4$}{$7$}
        \dijnode{v5}{dijclosed}{$v_5$}{$5$}
        \dijnode{v6}{dijopen}{$v_6$}{$7$}
        \dijnode{v7}{dijopen}{$v_7$}{$12$}
        \dijnode{v8}{dijopen}{$v_8$}{$10$}
        \dijnode{v9}{dijopen}{$v_9$}{$11$}
        \dijnode{v10}{dijopen}{$v_{10}$}{$17$}
    }
    \end{frame}

    \begin{frame}{Dijkstrův algoritmus}{Příklad}
    \centering
    \DijkstraState
    {Krok 6: uzavíráme $v_6$}
    {$v_9:9,\ v_8:10,\ v_7:10,\ v_{10}:14$}
    {
        \draw[dijtreeedge] (v1)--(v3);
        \draw[dijtreeedge] (v3)--(v2);
        \draw[dijtreeedge] (v3)--(v5);
        \draw[dijtreeedge] (v5)--(v4);
        \draw[dijtreeedge] (v5)--(v6);
    }
    {
        \dijnode{v1}{dijclosed}{$v_1$}{$0$}
        \dijnode{v2}{dijclosed}{$v_2$}{$3$}
        \dijnode{v3}{dijclosed}{$v_3$}{$2$}
        \dijnode{v4}{dijclosed}{$v_4$}{$7$}
        \dijnode{v5}{dijclosed}{$v_5$}{$5$}
        \dijnode{v6}{dijclosed}{$v_6$}{$7$}
        \dijnode{v7}{dijopen}{$v_7$}{$10$}
        \dijnode{v8}{dijopen}{$v_8$}{$10$}
        \dijnode{v9}{dijopen}{$v_9$}{$9$}
        \dijnode{v10}{dijopen}{$v_{10}$}{$14$}
    }
    \end{frame}

    \begin{frame}{Dijkstrův algoritmus}{Příklad}
    \centering
    \DijkstraState
    {Krok 7: uzavíráme $v_9$}
    {$v_7:10,\ v_8:10,\ v_{10}:11$}
    {
        \draw[dijtreeedge] (v1)--(v3);
        \draw[dijtreeedge] (v3)--(v2);
        \draw[dijtreeedge] (v3)--(v5);
        \draw[dijtreeedge] (v5)--(v4);
        \draw[dijtreeedge] (v5)--(v6);
        \draw[dijtreeedge] (v6)--(v9);
    }
    {
        \dijnode{v1}{dijclosed}{$v_1$}{$0$}
        \dijnode{v2}{dijclosed}{$v_2$}{$3$}
        \dijnode{v3}{dijclosed}{$v_3$}{$2$}
        \dijnode{v4}{dijclosed}{$v_4$}{$7$}
        \dijnode{v5}{dijclosed}{$v_5$}{$5$}
        \dijnode{v6}{dijclosed}{$v_6$}{$7$}
        \dijnode{v7}{dijopen}{$v_7$}{$10$}
        \dijnode{v8}{dijopen}{$v_8$}{$10$}
        \dijnode{v9}{dijclosed}{$v_9$}{$9$}
        \dijnode{v10}{dijopen}{$v_{10}$}{$11$}
    }
    \end{frame}

    \begin{frame}{Dijkstrův algoritmus}{Příklad}
    \centering
    \DijkstraState
    {Krok 8: uzavíráme $v_7$}
    {$v_8:10,\ v_{10}:11$}
    {
        \draw[dijtreeedge] (v1)--(v3);
        \draw[dijtreeedge] (v3)--(v2);
        \draw[dijtreeedge] (v3)--(v5);
        \draw[dijtreeedge] (v5)--(v4);
        \draw[dijtreeedge] (v5)--(v6);
        \draw[dijtreeedge] (v6)--(v9);
        \draw[dijtreeedge] (v6)--(v7);
    }
    {
        \dijnode{v1}{dijclosed}{$v_1$}{$0$}
        \dijnode{v2}{dijclosed}{$v_2$}{$3$}
        \dijnode{v3}{dijclosed}{$v_3$}{$2$}
        \dijnode{v4}{dijclosed}{$v_4$}{$7$}
        \dijnode{v5}{dijclosed}{$v_5$}{$5$}
        \dijnode{v6}{dijclosed}{$v_6$}{$7$}
        \dijnode{v7}{dijclosed}{$v_7$}{$10$}
        \dijnode{v8}{dijopen}{$v_8$}{$10$}
        \dijnode{v9}{dijclosed}{$v_9$}{$9$}
        \dijnode{v10}{dijopen}{$v_{10}$}{$11$}
    }
    \end{frame}

    \begin{frame}{Dijkstrův algoritmus}{Příklad}
    \centering
    \DijkstraState
    {Krok 9: uzavíráme $v_8$}
    {$v_{10}:11$}
    {
        \draw[dijtreeedge] (v1)--(v3);
        \draw[dijtreeedge] (v3)--(v2);
        \draw[dijtreeedge] (v3)--(v5);
        \draw[dijtreeedge] (v5)--(v4);
        \draw[dijtreeedge] (v5)--(v6);
        \draw[dijtreeedge] (v6)--(v9);
        \draw[dijtreeedge] (v6)--(v7);
        \draw[dijtreeedge] (v5)--(v8);
    }
    {
        \dijnode{v1}{dijclosed}{$v_1$}{$0$}
        \dijnode{v2}{dijclosed}{$v_2$}{$3$}
        \dijnode{v3}{dijclosed}{$v_3$}{$2$}
        \dijnode{v4}{dijclosed}{$v_4$}{$7$}
        \dijnode{v5}{dijclosed}{$v_5$}{$5$}
        \dijnode{v6}{dijclosed}{$v_6$}{$7$}
        \dijnode{v7}{dijclosed}{$v_7$}{$10$}
        \dijnode{v8}{dijclosed}{$v_8$}{$10$}
        \dijnode{v9}{dijclosed}{$v_9$}{$9$}
        \dijnode{v10}{dijopen}{$v_{10}$}{$11$}
    }
    \end{frame}

    \begin{frame}{Dijkstrův algoritmus}{Příklad}
    \centering
    \DijkstraState
    {Krok 10: uzavíráme $v_{10}$ -- konec}
    {prázdná}
    {
        \draw[dijtreeedge] (v1)--(v3);
        \draw[dijtreeedge] (v3)--(v2);
        \draw[dijtreeedge] (v3)--(v5);
        \draw[dijtreeedge] (v5)--(v4);
        \draw[dijtreeedge] (v5)--(v6);
        \draw[dijtreeedge] (v6)--(v9);
        \draw[dijtreeedge] (v6)--(v7);
        \draw[dijtreeedge] (v5)--(v8);
        \draw[dijtreeedge] (v9)--(v10);
    }
    {
        \dijnode{v1}{dijclosed}{$v_1$}{$0$}
        \dijnode{v2}{dijclosed}{$v_2$}{$3$}
        \dijnode{v3}{dijclosed}{$v_3$}{$2$}
        \dijnode{v4}{dijclosed}{$v_4$}{$7$}
        \dijnode{v5}{dijclosed}{$v_5$}{$5$}
        \dijnode{v6}{dijclosed}{$v_6$}{$7$}
        \dijnode{v7}{dijclosed}{$v_7$}{$10$}
        \dijnode{v8}{dijclosed}{$v_8$}{$10$}
        \dijnode{v9}{dijclosed}{$v_9$}{$9$}
        \dijnode{v10}{dijclosed}{$v_{10}$}{$11$}
    }
    \end{frame}

    \subsection{Správnost Dijkstrova algoritmu}

    \begin{frame}{Proč to funguje?}
        \visible<2->{\begin{propblock}
            Když Dijkstrův algoritmus otevírá vrchol $v$, pak platí, že $D(v)=d(v,v_0)$.
        \end{propblock}}
        \begin{itemize}
            \visible<3->{\item Jinak řečeno, když si algoritmus vybere vrchol $v$, aby jej otevřel, jeho vzdálenost od počátku je \textbf{spočítaná správně}.}
            \visible<4->{\item Na tuto skutečnost můžeme nahlédnout sporem.}
        \end{itemize}
    \end{frame}

    \begin{frame}{Proč to funguje?}
        \visible<2->{\begin{talertframe}[\textwidth]
            \begin{center}
                Co by se stalo, kdyby algoritmus vybral vrchol $v$, jehož $D$ neodpovídá délce nejkratší cesty?
            \end{center}
        \end{talertframe}}
        \visible<3->{\begin{figure}
            \centering
            \begin{tikzpicture}[scale=0.7]

    % \tikzset{
    %     vertex/.style={
    %         draw,
    %         circle,
    %         thick,
    %         fill=white,
    %         minimum size=7mm,
    %         inner sep=0pt,
    %         font=\small
    %     }
    % }

    % vrcholy
    \coordinate (v)  at (0,1.8);
    \coordinate (w)  at (5.6,4.0);
    \coordinate (v0) at (7.5,-1.2);

    % hrana v--w
    \draw[thick] (v) -- (w);

    % červená cesta D(v)
    \draw[red!80!black, line width=0.6pt,dashed]
        (v)
        .. controls (1.3,2.2) and (1.8,1.2) ..
        (1.0,0.8)
        .. controls (0.2,0.4) and (0.5,-0.4) ..
        (1.7,-0.8)
        .. controls (2.8,-1.2) and (3.4,-0.3) ..
        (4.0,0.6)
        .. controls (4.4,1.2) and (4.7,0.8) ..
        (4.9,-0.6)
        .. controls (5.1,-1.6) and (5.7,-1.5) ..
        (6.4,-1.0)
        .. controls (6.9,-0.6) and (7.2,-0.7) ..
        (v0);

    % modrá cesta D(w)
    \draw[blue!80!black, line width=0.6pt,dashed]
        (w)
        .. controls (5.3,3.5) and (5.1,2.9) ..
        (5.8,2.9)
        .. controls (6.6,2.9) and (6.9,2.4) ..
        (7.4,2.1)
        .. controls (8.2,1.5) and (8.3,0.8) ..
        (7.7,0.3)
        .. controls (7.1,-0.1) and (6.2,0.0) ..
        (6.4,-0.4)
        .. controls (6.6,-0.8) and (7.5,-0.4) ..
        (8.1,-0.8)
        .. controls (8.5,-1.1) and (8.2,-1.5) ..
        (v0);

    % vrcholy s popisky uvnitř
    \node[vertex,fill=white] at (v) {$v$};
    \node[vertex,fill=white] at (w) {$s$};
    \node[vertex,fill=white] at (v0) {$v_0$};

    % popisky
    \node at (2.3,3.25) {$w(v,s)$};
    \node[color=red] at (1.0,-1.05) {$D(v)$};
    \node[color=blue] at (8.2,2.0) {$D(s)$};

\end{tikzpicture}
        \end{figure}}
        \begin{itemize}
            \visible<4->{\item $D(v)$ = spočítaná cesta do $v$.}
            \visible<5->{\item $D(s)+w(v,s)$ = nejkratší cesta do $v$ vedoucí přes souseda $s$.}
        \end{itemize}
    \end{frame}

    \begin{frame}{Proč to funguje?}
        \begin{itemize}
            \visible<2->{\item Protože $D(v)$ je větší než délka nejkratší cesty do $v$, platí
            \[D(v)>D(s)+w(v,s).\]}
            \visible<3->{\item Připomeňme si nyní, že váhy všech hran jsou nezáporné.}
            \visible<4->{\item Pak tedy platí:
            \[D(v)>D(s)+\underbrace{w(v,s)}_{\geqslant 0}\geqslant D(s).\]}
            \visible<5->{\item To znamená, že \textbf{$D(s)$ je nižší než $D(v)$}.}
            \visible<6->{\item \markdarkred{To je ale spor}, neboť algoritmus by si tedy \textbf{nevybral} vrchol $v$.}
        \end{itemize}
    \end{frame}

    \begin{frame}{Nezápornost vah}
        \visible<2->{\begin{talertframe}
            \begin{center}
                Všimněte si, že právě díky \textbf{nezápornosti vah} hran grafu na vstupu algoritmus funnguje!
            \end{center}
        \end{talertframe}}
        \visible<3->{\begin{tnoticeframe}[\textwidth]
            \begin{center}
                Kdyby hrany mohly mít i záporné váhy, nejkratší sled mezi některými dvojicemi vrcholů vůbec nemusí existovat!
            \end{center}
        \end{tnoticeframe}}
        \visible<4->{\begin{figure}
            \centering
            \input{images/zaporny-cyklus.tex}
        \end{figure}}
    \end{frame}

    \begin{frame}{Nezápornost vah}
        \visible<2->{\begin{talertframe}
            \begin{center}
                Dále se také může stát, že hledání nejkratší cesty přejde na problém hledání \textbf{Hamiltonovské cesty} $\implies$ $\classNP$-úplný problém.
            \end{center}
        \end{talertframe}}
        \visible<3->{\begin{figure}
            \centering
            \begin{tikzpicture}[
    scale=1,
    vertex/.style={draw, circle, thick, minimum size=7mm, inner sep=0pt},
    diredge/.style={->, thick, >=stealth},
    bestedge/.style={->, very thick, draw=red!70!black, >=stealth},
    weight/.style={fill=white, inner sep=1pt, font=\small}
]

\node[vertex] (s)  at (0,0) {$s$};
\node[vertex] (v1) at (2,0) {$v_1$};
\node[vertex] (v2) at (4,0) {$v_2$};
\node[vertex] (v3) at (6,0) {$v_3$};
\node[vertex] (v4) at (8,0) {$v_4$};
\node[vertex] (t)  at (10,0) {$t$};

% Hamiltonovská cesta
\draw[bestedge] (s)  -- node[weight] {$-1$} (v1);
\draw[bestedge] (v1) -- node[weight] {$-1$} (v2);
\draw[bestedge] (v2) -- node[weight] {$-1$} (v3);
\draw[bestedge] (v3) -- node[weight] {$-1$} (v4);
\draw[bestedge] (v4) -- node[weight] {$-1$} (t);

% "kratší" lokální alternativy, které ale nejsou globálně nejlepší
\draw[diredge] (s) to[bend left=35] node[weight] {$0$} (v2);
\draw[diredge] (s) to[bend left=45] node[weight] {$0$} (v3);
\draw[diredge] (s) to[bend left=55] node[weight] {$0$} (t);

\draw[diredge] (v1) to[bend left=35] node[weight] {$0$} (v3);
\draw[diredge] (v2) to[bend left=35] node[weight] {$0$} (v4);
\draw[diredge] (v3) to[bend left=35] node[weight] {$0$} (t);

\node[align=center] at (5,-1.4)
{\large Nejkratší cesta z~$s$ do~$t$ je\\[1mm]
$s\to v_1\to v_2\to v_3\to v_4\to t$ s~cenou $-5$};

\end{tikzpicture}
        \end{figure}}
    \end{frame}

    \section{Časová a prostorová složitost}

    \titleframe{Složitost algoritmu}

    \subsection{Kam ukládat vrcholy?}

    \begin{frame}{Kam ukládat vrcholy?}
        \begin{itemize}
            \visible<2->{\item Zatím jsme vůbec nevyřešili, jak budeme ukládat hodnoty $D$.}
            \visible<3->{\item Jako první varianta se nabízí \markdarkred{obyčejný seznam/pole}.}
        \end{itemize}
        \visible<4->{\begin{tnoticeframe}[0.8\textwidth]
            Zkusíme určit \textbf{časovou} a \textbf{prostorovou složitost}.
        \end{tnoticeframe}}
    \end{frame}

    \subsection{Složitost s polem}

    \begin{frame}{Časová složitost s polem}
        \begin{itemize}
            \visible<2->{\item Opět si označíme $n=|V|$ a $m=|E|$.}
            \visible<3->{\item Inicializace potrvá $\bigO(2n)=\bigO(n)$, neboť zavádíme seznam stavů a vzdáleností \textbf{pro všechny vrcholy}.}
            \visible<4->{\item Každý vrchol uzavřeme \textbf{nejvýše jednou} $\implies$ vnější cyklus provede nejvýše $\bigO(n)$ iterací.}
            \visible<5->{\item Hledáním minimální hodnoty $D$ strávíme $\bigO(n)$ $\impliedby$ \textbf{hledáme minimum v poli}.}
            \visible<6->{\item Každou hranu při procházení sousedů projdeme \textbf{nejvýše dvakrát} $\implies$ to dává řádově $\bigO(2m)=\bigO(m)$ kroků.}
            \visible<7->{\item Aktualizaci hodnoty $D(v)$ libovolného vrcholu $v$ provedeme v konstantním čase.}
            \visible<8->{\item Dohromady:
            \[\bigO(\underbrace{n}_\text{inicializace}+\overbrace{n}^{\text{vnejší cyklus}}\cdot\overbrace{n}^{\text{minimum}}+\underbrace{m}_\text{sousedé})=\bigO(n^{2}+m).\]}
        \end{itemize}
    \end{frame}

    \begin{frame}{Prostorová složitost s polem}
        \begin{itemize}
            \visible<2->{\item Pokud budeme reprezentovat graf pomocí \textbf{seznamu sousedů}, spotřebujeme $\bigO(n+m)$ paměti.}
            \visible<3->{\item Dále již potřebujeme jen pole pro uchování stavů a vzdáleností.}
            \visible<4->{\item Celkově činí prostorová složitost $\bigO(n+m)$.}
        \end{itemize}
    \end{frame}

    \begin{frame}{Shrnutí}
        \visible<2->{S polem jsme dosáhli následujících výsledků:}
        \begin{itemize}
            \visible<3->{\item Časová složitost: $\bigO(n^{2}+m)$.}
            \visible<4->{\item Prostorová složitost $\bigO(n+m)$.}
        \end{itemize}
        \visible<5->{\begin{tnoticeframe}
            \begin{center}
                S prostorovou složitostí toho moc nezmůžeme, nicméně v případě časové složitosti \textbf{lze výsledek o něco zlepšit} $\implies$ máme přeci \markdarkred{binární haldu}!
            \end{center}
        \end{tnoticeframe}}
        \visible<6->{Ta se nám bude hodit v případě \textbf{výběru minima} z hodnot $D$.}
    \end{frame}

    \subsection{Složitost s binární haldou}

    \begin{frame}{Časová složitost s binární haldou}
        \begin{itemize}
            \item Z počátku spotřebujeme $\bigO(n)$ času na zavedení pole stavů.
            \item V případě haldy provedeme $n$-krát operaci $\textsc{Insert}$, tzn. spotřebujeme $\bigO(n\log{n})$ času.
            \item Při výběru vrcholu s nejnižším $D$ provádíme operaci $\textsc{ExtractMin}$. Celkově ji provedeme $n$-krát (vnější cyklus), tzn. spotřebujeme $\bigO(n\log{n})$ času.
            \item Při aktualizaci hodnoty $D(v)$ libovolného vrcholu provádíme operaci $\textsc{Decrease}$. Nejvýše jej provedeme $m$-krát (za každou hranu) $\implies$ spotřebujeme $\bigO(m\log{n})$.
            \item Dohromady:
            \begin{align*}
                &\bigO(\underbrace{n\log{n}}_\text{inicializace}+\overbrace{n}^\text{vnější cyklus}\cdot\overbrace{\log{n}}^{\textsc{ExtractMin}}+\underbrace{m}_\text{hrany}\cdot\underbrace{\log{n}}_{\textsc{Decrease}}) \\
                &=\bigO((n+m)\log{n}).
            \end{align*}
        \end{itemize}
    \end{frame}

    \subsection{Shrnutí a závěrečný rozbor}

    \begin{frame}{Shrnutí a závěrečný rozbor}
        \visible<2->{\begin{tinfoframe}
            Časová složitost \textbf{s polem} vs. \textbf{s haldou}:
            \begin{itemize}
                \item Dijkstrův algoritmus s polem: $\bigO(n^{2}+m)$
                \item Dijkstrův algoritmus s binární haldou: $\bigO((n+m)\log{n})$
            \end{itemize}
        \end{tinfoframe}}
        \visible<3->{Nyní se nybízí dvojice případů:}
        \begin{itemize}
            \visible<4->{\item Graf $G$ je řídký, tj. má \markdarkred{nejvýše lineárně mnoho} hran vzhledem k počtu vrcholů, tj. $m=\bigO(n)$. Pak}
            \begin{itemize}
                \visible<5->{\item S polem: $\bigO(n^{2}+m)=\bigO(n^{2}+n)=\bigO(n^{2})$.}
                \visible<6->{\item S haldou: $\bigO((n+m)\log{n})=\bigO(2n\log{n})=\bigO(n\log{n})$}
            \end{itemize}
            \visible<7->{\item Graf $G$ je hustý, tj. má \markdarkred{kvadraticky mnoho} hran vzhledem k počtu vrcholů, tj. $m=\bigO(n^{2})$. Pak}
            \begin{itemize}
                \visible<8->{\item S polem: $\bigO(n^{2}+n^{2})=\bigO(2n^{2})=\bigO(n^{2})$.}
                \visible<9->{\item S haldou: $\bigO((n+n^{2})\log{n})=\bigO(n\log{n}+n^{2}\log{n})=\bigO(n^{2}\log{n})$.}
            \end{itemize}
        \end{itemize}
    \end{frame}

    \begin{frame}{Shrnutí a závěrečný rozbor}
        \begin{timportantframe}[\textwidth]
            \begin{itemize}
                \item Pro \textbf{řídké} grafy vychází asymptotická složitost lépe v případě \markdarkred{implementace s binární haldou}.
                \item Pro \textbf{husté} grafy vychází lépe \markdarkred{implementace s polem}.
            \end{itemize}
        \end{timportantframe}
    \end{frame}

    \section{Zdroje}

    \begin{frame}{Zdroje}
        \begin{itemize}
            \item MAREŠ, Martin a~VALLA, Tomáš. \textit{Průvodce labyrintem algoritmů.} Druhé vydání. CZ.NIC. Praha: CZ.NIC, z.s.p.o., 2022. ISBN 978-80-88168-63-8.
            \item CORMEN, Thomas H.; LEISERSON, Charles Eric; RIVEST, Ronald L. a~STEIN, Clifford. \textit{Introduction to algorithms.} Fourth edition. Cambridge: The MIT Press, 2022. ISBN 978-0-262-04630-5.
            \item MATOUŠEK, Jiří a NEŠETŘIL, Jaroslav. \textit{Kapitoly z diskrétní matematiky.} 4., upr. a dopl. vyd. V Praze: Karolinum, 2009. ISBN 978-80-246-1740-4.
        \end{itemize}
    \end{frame}

\end{document}